% Source: https://artofproblemsolving.com/community/c3h1210527_challenging_question
% Stable identity: aops-topic-1210527-q01
The triangle returns to the same position after any rotation by a multiple of $120^\circ$, so only the total angle modulo $120^\circ$ matters. After $n$ steps the total is
\[s_n=3+9+\cdots+3^n=\dfrac{3^{n+1}-3}{2}.\]
Modulo $120$: $s_0=0$, $s_1=3$, $s_2=12$, $s_3=39$, $s_4=120\equiv0$. In general $s_{n+4}-s_n=3^{n+1}(1+3+9+27)=40\cdot3^{n+1}\equiv0\pmod{120}$ (as $3\mid3^{n+1}$). So the totals repeat with period $4$: $0,3,12,39$.

\begin{asy}
usepackage("amsmath");
size(3cm);
defaultpen(fontsize(10pt));
pair O=(0,0); for(real t: new real[]{0,3,12,39}) draw(rotate(-t)*(dir(90)--dir(210)--dir(330)--cycle),t==0?blue+linewidth(1):gray); dot(O);
label("\mbox{}",(min(currentpicture,user=true)+max(currentpicture,user=true))/2,black); // keeps the text after the figure black when LaTeX typesets the labels inline
\end{asy}

The triangle occupies $\boxed{4}$ different positions.
